% Generated draft: requires mathematical and source review.
\documentclass[11pt]{article}
\usepackage[margin=25mm]{geometry}
\usepackage{fontspec}
\setmainfont{Latin Modern Roman}
\usepackage{amsmath,amssymb,mathtools,unicode-math}
\setmathfont{Latin Modern Math}
\usepackage{enumitem,needspace,xurl,hyperref}
\setlength{\parindent}{0pt}
\setlength{\parskip}{5pt}
\setlength{\emergencystretch}{3em}
\newcommand{\N}{\mathbb{N}}
\newcommand{\Z}{\mathbb{Z}}
\newcommand{\Q}{\mathbb{Q}}
\newcommand{\R}{\mathbb{R}}
\newcommand{\C}{\mathbb{C}}
\newcommand{\F}{\mathbb{F}}
\newcommand{\A}{\mathbb{A}}
\newcommand{\PP}{\mathbb{P}}
\newcommand{\E}{\mathbb{E}}
\newcommand{\eps}{\varepsilon}
\newcommand{\dd}{\,\mathrm{d}}
\begin{document}
\section*{UnsolvedMath ID: 6800022}
\textbf{Source problem code:} catalogue entry
\subsection{Definitions and mathematical statement}
Use the usual setting of a smooth embedded minimal surface \(\Sigma\subset\mathbb R^3\), with its induced metric and zero mean curvature. Completeness refers to this intrinsic metric. The inclusion is proper if its inverse image of every compact subset of \(\mathbb R^3\) is compact in \(\Sigma\).

(a) The Calabi--Yau question asks whether every complete embedded minimal surface in \(\mathbb R^3\) is properly embedded. No finite-genus or finite-total-curvature condition is present in this part. The excerpt does not explicitly give boundary or connectedness conventions; the standard connected boundaryless surface interpretation requires confirmation if alternatives are intended.

(b) The Hoffman--Meeks conjecture concerns complete embedded minimal surfaces of finite genus \(g\), with \(k\) ends and finite total curvature. With Gaussian curvature \(K_\Sigma\), the usual minimal-surface finiteness convention is
\[
 \int_\Sigma |K_\Sigma|\,\dd A<\infty.
\]
An end is a topological end, represented by a coherent component outside compact subsets. The proposed inequality is
\[
 k\le g+2.
\]
All the hypotheses in this part apply together; they are not transferred to part (a). The section heading about constant mean curvature in homogeneous spaces does not enlarge these two explicit questions beyond minimal surfaces in \(\mathbb R^3\).

The supplied summary labels (a) negatively solved in general and attributes counterexamples to ``Meeks/Nakamori/Rosenberg'' and 2017--2019 constructions, but gives no identifiable publication or precise example. This unsupported report is not certified as a resolution of the exact embedded Euclidean question; immersed examples or examples in other settings would not by themselves resolve it. The summary separately calls (b) open as of 2020 and cites a Meeks--P\'erez--Ros problem survey without full details. The registry still has an open label for the combined record. These dated, mixed status claims are retained explicitly without an independent modern literature assessment.
\subsection{Short English statement}
Does completeness force properness for embedded minimal surfaces in \(\mathbb R^3\)? Separately, must a complete embedded finite-total-curvature surface with genus \(g\) have at most \(g+2\) ends? The summary's alleged negative resolution of the first part lacks a precise supporting reference.
\subsection{Sources}
\begin{itemize}
\item Local registry: \url{data/unsolvedmath/problems.json}, permanent UnsolvedMath ID \(6800022\), source code catalogue entry; supplied statement, background, and research summary, updated 2026-08-06. Proposed by Joaquín P\'erez, University of Granada, Spain, 2018.
\item Morgan and Pansu, \emph{A List of Open Problems in Differential Geometry} (2018), Question environment 22, ``Constant mean curvature in homogeneous 3-manifolds'': \url{https://www.imo.universite-paris-saclay.fr/~pansu/problems_MTDG.tex}. Supplied source-TeX extraction accessed 2026-07-29.
\item The summary mentions ``Meeks/Nakamori/Rosenberg'' counterexamples and 2017--2019 constructions without full citations, and a Meeks--P\'erez--Ros problem survey for its 2020 Hoffman--Meeks status report, again without exact publication details.
\end{itemize}
\end{document}
